%%%%%%%%%%%%%%%%%%%%%%%%%%%%%%%%%%%%%%%%%%%%%%%%%%%%%%%%%%%%
%%  This Beamer template was created by Cameron Bracken.
%%  Anyone can freely use or modify it for any purpose
%%  without attribution.
%%
%%  Last Modified: January 9, 2009
%%

\documentclass[xcolor=x11names,compress]{beamer}

%% General document %%%%%%%%%%%%%%%%%%%%%%%%%%%%%%%%%%
\usepackage{graphicx}
\usepackage{tikz}
\usetikzlibrary{decorations.fractals}
%%%%%%%%%%%%%%%%%%%%%%%%%%%%%%%%%%%%%%%%%%%%%%%%%%%%%%

\setbeamertemplate{footline}[frame number]

%% Beamer Layout %%%%%%%%%%%%%%%%%%%%%%%%%%%%%%%%%%
\useoutertheme[subsection=false,shadow]{miniframes}
\useinnertheme{default}
\usefonttheme{serif}
\usepackage{palatino}
%\setbeamertemplate{footline}[frame number]

\usepackage{float}
\usepackage{gensymb}
\usepackage{amssymb}
\usepackage{amsmath}
\usepackage{mathtools}
\usepackage{ellipsis}
\usepackage{textcomp}
\usepackage{cite}
\usepackage{graphicx}%[pdftex]
\usepackage{subfigure}
\usepackage{epsfig}
\usepackage{caption}
\usepackage{algorithm}
%\usepackage[linesnumbered,ruled,vlined]{algorithm2e}
\usepackage{algpseudocode}
\usepackage{algorithmicx}
\usepackage{color}
\usepackage{epstopdf}
\usepackage{psfrag}
\usepackage{grffile}
\DeclareMathOperator*{\Max}{Max}
\usepackage{setspace}
\usepackage{amsthm}
\newtheorem{thm}{Theorem}
\theoremstyle{definition}
\newtheorem{dfn}{Definition}
\theoremstyle{remark}
%\newtheorem{note}{Note}[section]
\theoremstyle{plain}
\newtheorem{lem}{Lemma}
\newtheorem{proposition}{Proposition}
\newtheorem{crl}{Corollary}
\newtheorem{assump}{Assumption}
\usepackage{xcolor,colortbl}
\usepackage{pifont}
\usepackage{xcolor}
\definecolor{bostonuniversityred}{rgb}{0.8, 0.0, 0.0}
\definecolor{burgundy}{rgb}{0.5, 0.0, 0.13}

\usepackage{hhline}
\usepackage{multirow}
\DeclareMathOperator*{\argmax}{arg\,max}

\setbeamerfont{title like}{shape=\scshape}
\setbeamerfont{frametitle}{shape=\scshape}

\setbeamercolor*{lower separation line head}{bg=DeepSkyBlue4} 
\setbeamercolor*{normal text}{fg=black,bg=white} 
\setbeamercolor*{alerted text}{fg=red} 
\setbeamercolor*{example text}{fg=black} 
\setbeamercolor*{structure}{fg=black} 
 
\setbeamercolor*{palette tertiary}{fg=black,bg=black!10} 
\setbeamercolor*{palette quaternary}{fg=black,bg=black!10} 

\renewcommand{\(}{\begin{columns}}
\renewcommand{\)}{\end{columns}}
\newcommand{\<}[1]{\begin{column}{#1}}
\renewcommand{\>}{\end{column}}
%%%%%%%%%%%%%%%%%%%%%%%%%%%%%%%%%%%%%%%%%%%%%%%%%%




\begin{document}


%%%%%%%%%%%%%%%%%%%%%%%%%%%%%%%%%%%%%%%%%%%%%%%%%%%%%%
%%%%%%%%%%%%%%%%%%%%%%%%%%%%%%%%%%%%%%%%%%%%%%%%%%%%%%

\begin{frame}
\title{Technical Report Writing using \LaTeX}
%\subtitle{SUBTITLE}
%\author{
%	Cameron Bracken\\
%	{\it Humboldt State University}\\
%}

\author[\vspace*{-.5cm} Department of Electronics \& Communication]{ \footnotesize{Presentation\\} \footnotesize{by\\} \normalsize{Dr. Shubhankar Majumdar} }
\vspace*{-.6cm}
\date{ \vspace*{-1.2cm}
\begin{figure}[!ht]
\centering
\includegraphics[scale=0.3]{logo}
\end{figure}
\vspace{-.4cm}
\scriptsize{Department of Electronics \& Communication\\}
\scriptsize{National Institute of Technology Meghalaya\\} 
\today 
}

\titlepage
\end{frame}


%\begin{frame}
%\vspace*{-2.3cm} 
%\titlepage % Print the title page as the first slide
%\end{frame}

\begin{frame}
\frametitle{{\color{DeepSkyBlue4}Outline}} 
\small{
\begin{itemize}


%
\item Why \LaTeX?
\item Introduction to \LaTeX
\item Documents Structure of \LaTeX
\item Few Basics (EXAMPLE: 1 (Simple Document) \& EXAMPLE: 2 (Two Column Document of IEEEtran class))
\item Mathematical Equation in \LaTeX
\item Mathematical Equation in \LaTeX (EXAMPLE: 3)
\item Creating Tables in \LaTeX (EXAMPLE: 4 \& EXAMPLE: 5 )
\item Importing Figures in \LaTeX (EXAMPLE: 6 \& EXAMPLE: 7 )
\item Reference \& Citation (EXAMPLE: 8)
\item How to write proofs of theorems (EXAMPLE: 9)
\item How to write Algorithms (EXAMPLE: 10)

\end{itemize}
}
\end{frame}

%----------------------------------------------------------------------------------------
%	PRESENTATION SLIDES
%----------------------------------------------------------------------------------------

%------------------------------------------------
\begin{frame}
\frametitle{{\color{DeepSkyBlue4}Why \LaTeX?}} 
\begin{enumerate}
\item 	It's \textbf{free} and \textbf{portable}.
\item  You can use the editor of your choice (Even MS Word). 
\item Including \textbf{mathematical expressions} in LaTeX involves typing a few appropriate characters. By contrast, including mathematics in Word requires Equation Editor, a cumbersome and slow graphical user interface 
\item  \textbf{Style changes are neater} in \LaTeX. \textbf{Style files for many periodicals exist}.  
\item Almost all mathematical and scientific \textbf{notations} are easily achievable in LaTeX.
\item \textbf{Repetitive tasks} can more easily be \textbf{automated}. 
\item Don't bother about the \textbf{format}, concentrate on the \textbf{content}. 
\end{enumerate}
\end{frame}


\end{document}
